% Gurvir Singh — one-page CV, aimed at AI engineering internships
%
% Build (from the project root, with a LaTeX distribution such as MiKTeX):
%   pdflatex -output-directory=cv/build cv/gurvir-singh-cv.tex
% then copy cv/build/gurvir-singh-cv.pdf to public/files/gurvir-singh-cv.pdf
%
% Typography follows a classic LaTeX academic CV: Computer Modern (T1 / cm-super),
% 11pt body, bold small-caps name and section headings, full-width section rules.

\documentclass[11pt]{article}

\usepackage[T1]{fontenc}
% Use the scalable Type 1 cm-super fonts (some installations do not map them by default).
\pdfmapfile{+cm-super-t1.map}
\usepackage[a4paper, left=0.75in, right=0.75in, top=0.6in, bottom=0.5in]{geometry}
\usepackage{microtype}
\usepackage{enumitem}
\usepackage{titlesec}
\usepackage[hidelinks,
  pdftitle={Gurvir Singh - Curriculum Vitae},
  pdfauthor={Gurvir Singh},
  pdfsubject={AI engineering internship},
  pdfkeywords={AI engineering, LLM, machine learning, Python}]{hyperref}

\pagestyle{empty}
\setlength{\parindent}{0pt}
\setlength{\parskip}{0pt}

% ── Section headings: large bold small caps over a full-width rule ───────────
\titleformat{\section}{\Large\bfseries\scshape\raggedright}{}{0pt}{}[{\vspace{-2pt}\titlerule[1.3pt]}]
\titlespacing*{\section}{0pt}{6pt}{3pt}

% ── Entries ─────────────────────────────────────────────────────────────────
% \entry{bold left}{bold right}
\newcommand{\entry}[2]{\hspace*{6pt}\textbf{#1}\hfill\textbf{#2}\par}
% \detail{text}{right-aligned text}, indented under an entry
\newcommand{\detail}[2]{\hspace*{27pt}#1\hfill#2\par}
% \project{name}{url}{note}{stack}
\newcommand{\project}[4]{%
  \hspace*{6pt}\textbf{\href{#2}{#1}}#3\hfill\textit{#4}\par}
% \projectnolink{name}{note}{stack}, for work without a public repository
\newcommand{\projectnolink}[3]{%
  \hspace*{6pt}\textbf{#1}#2\hfill\textit{#3}\par}

% Bullet lists
\setlist[itemize]{label=$\bullet$, leftmargin=34pt, labelsep=6pt,
  topsep=2pt, itemsep=1pt, parsep=0pt, partopsep=0pt}
\newenvironment{projectitems}
  {\begin{itemize}[topsep=1pt, itemsep=0.5pt]}
  {\end{itemize}\vspace{3pt}}

\begin{document}

% ── Name and contact ────────────────────────────────────────────────────────
\begin{center}
  {\LARGE\bfseries\scshape Gurvir Singh}\\[2pt]
  Third-Year B.Tech Student \textbullet\ Seeking an AI Engineering Internship
\end{center}
\vspace{-4pt}

\noindent
\hspace*{6pt}%
\begin{tabular}[t]{@{}l@{}}
  Computer Science \& Engineering\\
  (Artificial Intelligence \& Data Science)\\
  GNA University, Punjab, India
\end{tabular}%
\hfill
\begin{tabular}[t]{@{}r@{}}
  \href{mailto:gurvir.singh.panesar@gmail.com}{gurvir.singh.panesar@gmail.com} \textbullet\ +91 9592592881\\
  \href{https://github.com/GurvirSinghH}{github.com/GurvirSinghH}\\
  \href{https://gurvirsingh.in}{gurvirsingh.in}
\end{tabular}%
\hspace*{3pt}

% ── Profile ─────────────────────────────────────────────────────────────────
\section{Profile}
\hspace*{6pt}\parbox{\dimexpr\linewidth-6pt}{%
  Computer Science (AI \& Data Science) student who builds working AI tools end to end in Python: an
  LLM assistant that writes code, runs it and feeds errors back to the model to fix; an ML app that
  flags anomalies in system logs; and a crawler that turns a website into a retrieval-ready dataset.}

% ── Projects ────────────────────────────────────────────────────────────────
\section{Projects}

\project{Neural Forge (AI Blender Assistant)}{https://github.com/GurvirSinghH/neural-forge}%
  { (v1.0.0)}{Python, Blender Python API, LLM APIs}
\begin{projectitems}
  \item Blender add-on that turns natural-language prompts into Blender Python scripts, with one
        client for Gemini, OpenAI, local Ollama models and any OpenAI-compatible endpoint.
  \item Adds scene context (objects, selection) to the prompt, runs the generated code and sends errors
        back to the model to fix them: an agent-style generate, run and repair loop.
\end{projectitems}

\project{AI Log Intelligence Platform}{https://github.com/GurvirSinghH/AI-Log-Intelligence-Platform}%
  { (in development)}{Python, Streamlit, Pandas, Scikit-learn, Plotly}
\begin{projectitems}
  \item Streamlit app that detects and parses Linux auth/syslog and Apache access/error logs, with summary
        statistics, keyword search, Plotly charts and an optional Gemini-drafted incident report.
  \item Label-encoded process, module and host fields plus PID, message length and hour feed an
        Isolation Forest anomaly detector; K-Means then clusters the flagged entries.
\end{projectitems}

\project{College Knowledge Base Crawler}{https://github.com/GurvirSinghH/college-rag-assistant}%
  {}{Python, Requests, BeautifulSoup, lxml, SQLite}
\begin{projectitems}
  \item Multi-threaded crawler (worker pool and task queue) with robots.txt compliance, a shared rate
        limit, URL normalization, and retries with exponential backoff that respect Retry-After.
  \item Extracts page content and metadata to JSON, downloads linked documents and records crawl
        state in SQLite, as the data layer for a RAG pipeline (retrieval step not built yet).
\end{projectitems}

\projectnolink{Home Energy Monitor}{ (in progress)}{ESP32, PZEM-004T, Fedora Linux}
\begin{projectitems}
  \item An ESP32 with a PZEM-004T meter sends household power readings to my self-hosted Fedora server
        (message broker and logger); a simulated dry run works. Next: recognise which appliance is
        running from its power signature and estimate the monthly bill.
\end{projectitems}

\hspace*{6pt}\parbox{\dimexpr\linewidth-6pt}{%
  \textit{Also built:} \href{https://github.com/GurvirSinghH/CodeEditor}{CodeEditor} (Java Swing editor,
  MySQL over JDBC), \href{https://github.com/GurvirSinghH/csv-cleaner}{CSV Cleaner} (Python data cleaning),
  \href{https://github.com/GurvirSinghH/MLlab}{ML Lab} (Scikit-learn notebooks).}
\vspace{-2pt}

% ── Technical skills ────────────────────────────────────────────────────────
\section{Technical Skills}
\hspace*{6pt}%
\begin{tabular}{@{}l@{\hspace{12pt}}p{0.82\linewidth}@{}}
  \textit{Languages}       & Python (primary), C++, C, Java, HTML, CSS\\
  \textit{AI and LLMs}     & Gemini, OpenAI, Ollama and OpenAI-compatible APIs; prompt and context design;
                             LLM code generation with error feedback\\
  \textit{ML and data}     & Scikit-learn (Isolation Forest, K-Means, decision trees, regression), NumPy,
                             Pandas, Matplotlib, Seaborn, Plotly\\
  \textit{Tools}           & Streamlit, Git/GitHub, Jupyter, SQLite, MySQL, Linux (Fedora), SSH, Tailscale\\
  \textit{Fundamentals}    & Data structures and algorithms (practised in Python and C++)
\end{tabular}

% ── Education ───────────────────────────────────────────────────────────────
\section{Education}
\entry{GNA University}{Expected 2028}
\detail{B.Tech, Computer Science \& Engineering (AI \& Data Science), third year}{CGPA: 8.13/10}
\detail{Coursework: Data Visualization, Software Engineering, Theory of Computation, Computer Graphics}{}

% ── Currently learning ──────────────────────────────────────────────────────
\section{Currently Learning}
\begin{itemize}
  \item \textbf{AI engineering:} working through
        \href{https://github.com/rohitg00/ai-engineering-from-scratch}{AI Engineering from Scratch},
        an open-source course from math and ML foundations through deep learning to LLMs, RAG and agents.
  \item \textbf{Physics-informed ML:} building known physics into model training, e.g.\ via the loss function.
\end{itemize}

\end{document}
